\section{Research Experience}
\resumeSubHeading{Stanford University \& National University of Singapore}{Jun 2025 -- Sep 2026}
  {Research Assistant, Stanford Vision and Learning Lab and Adaptive Computing Lab}{}
Advised by Prof. \jiajun, Prof. \david, and Prof. \weiyu.
\begin{itemize}
  \item Led Planning in Abstract Worlds (PAW) as an independent researcher. We learned a natural-language abstract world model from unlabelled egocentric videos and used it to \textbf{reconstruct state--action supervision}. The same world model \textbf{conditions the planner} during supervised fine-tuning and inference, enabling \textbf{planning beyond demonstrated procedures} by recombining familiar actions to reach a given goal.
  \item On held-out CALVIN long-horizon tasks with goal-only instructions, PAW solved \textbf{57/60 (95\%)} versus 42/60 (70\%) for the fine-tuned Qwen3-VL planner baseline, a gain of \textbf{25 percentage points}.
  \item Against the same baseline, PAW raised the mean LLM-judged score from 4.43 to \textbf{4.75/5} on real-world test plans. Across three complex, long-horizon kitchen task domains, intervention-free completion rose from 24\% to \textbf{78\%} for new action sequences and from 11\% to \textbf{60\%} for new strategies.
\end{itemize}
\newpage
\resumeSubHeading{National University of Singapore}{Sep 2023 -- May 2024}
  {Undergraduate Research Assistant, GLOW.AI}{}
Advised by Prof. \href{https://www.comp.nus.edu.sg/cs/people/lowkh/}{Bryan Kian Hsiang Low}.
\begin{itemize}
  \item Designed and evaluated language-model planners for autonomous driving and household manipulation. My honours thesis examined how these planners choose actions under uncertainty.
\end{itemize}
\entryspace
\resumeSubHeading{University of Illinois Urbana-Champaign}{Oct 2023 -- Feb 2024}
  {Undergraduate Research Assistant, Robotic Perception, Interaction, and Learning Lab}{}
Advised by Prof. \yunzhu.
\begin{itemize}
  \item Helped develop RoboEXP, which uses interactive exploration to construct an action-conditioned scene graph. Supported deployment, evaluation, and visualization in OmniGibson and SAPIEN.
  \item Across five real-world exploration tasks, RoboEXP reconstructed the exact scene graph in \textbf{37/50 trials (74\%)}, versus 6/50 (12\%) for GPT with human-executed actions, a gain of \textbf{62 percentage points}.
\end{itemize}
\entryspace
\resumeSubHeading{Stanford University}{Apr 2023 -- Sep 2023}
  {Undergraduate Visiting Research Intern, Stanford Vision and Learning Lab}{}
Advised by Prof. \jiajun\ and Prof. \yunzhu.
\begin{itemize}
  \item Co-developed D\textsuperscript{3}Fields for robotic rearrangement from goal images. Built household manipulation tasks in OmniGibson and evaluated the representation in simulation and on a physical robot.
  \item Built the 3D representation from four RGB-D views in \textbf{0.166 seconds} on an A6000 GPU, versus 88.379 seconds for F3RM (about $532\times$ faster). On the Dense Object Nets dataset, D\textsuperscript{3}Fields achieved comparable correspondence precision and recall with \textbf{4 views instead of 50}.
\end{itemize}
\entryspace
\resumeSubHeading{National University of Singapore}{Jan 2021 -- Aug 2021}
  {Undergraduate Research Assistant, Adaptive Computing Lab}{}
Advised by Prof. \david.
\begin{itemize}
  \item Helped design and test a long-horizon visual navigation system for a Boston Dynamics Spot robot. Collected imitation-learning data on elevator use and navigation around pedestrians and obstacles.
\end{itemize}
